\part{Project Management and Team Collaboration}

\chapter{Project Management and Team Collaboration}

This section explains how the team intended to turn architecture into working software. The assignment
brief is explicit that repository organization, ownership, code review, and risk management are part of
the assessed design, not administrative afterthoughts.

\section{Team Structure and Ownership}

The team structure follows service and concern boundaries rather than assigning work randomly.

\begin{itemize}
  \item \textbf{Tabina} owned architecture, infrastructure, repository organization, timeline, and final
  consistency checks.
  \item \textbf{Noel} owned the backend, API contracts, auth/authorization design, and realtime feasibility.
  \item \textbf{Nurtore} owned AI behavior, cost/privacy considerations, and the data model.
  \item \textbf{Ming Ming} owned frontend UX, collaboration flows, and API usability from the client side.
\end{itemize}

Features spanning multiple owners were coordinated through contract-first discussions. For example,
realtime collaboration required alignment between frontend state behavior, backend role rules, and the
separate realtime service; AI assistance required alignment between UI review behavior, backend async job
contracts, AI provider behavior, versioning, and auditability.

When team members disagreed on a technical decision, the group resolved the disagreement by asking which
architectural driver mattered most for the feature in question. For example, demo reliability and
maintainability justified a deterministic stub AI provider and SQLite in the delivered MVP, even though
the longer-term architecture could target a fuller identity provider and larger-scale database deployment.

\section{Development Workflow}

\subsection{Branching Strategy}

The intended workflow is feature-branch based with protected \texttt{main}. Small, purpose-specific
branches reduce merge risk and keep code review focused. A representative naming convention is:

\begin{itemize}
  \item \texttt{backend/<name>-<feature>}
  \item \texttt{frontend/<name>-<feature>}
  \item \texttt{realtime/<name>-<feature>}
  \item \texttt{mvp/<milestone>} for final integration work
\end{itemize}

The merge policy is PR-first: no direct pushes to \texttt{main}. Every feature branch is expected to leave
the codebase runnable and testable.

Representative PR approval criteria were:

\begin{itemize}
  \item the change compiles or runs locally,
  \item the touched service tests pass or a manual verification note is included,
  \item API contracts and docs stay aligned,
  \item the commit history remains small, purposeful, and reviewable.
\end{itemize}

\subsection{Code Review Process}

The team used small commits and service-scoped reviews as the default. Review criteria were:

\begin{enumerate}
  \item does the change preserve the documented contract,
  \item is server-side authorization still enforced,
  \item does the feature have enough verification (tests or a justified manual check),
  \item does the change reduce or increase coupling,
  \item can a teammate understand the code quickly.
\end{enumerate}

Backend and realtime changes were reviewed especially carefully when they affected collaboration,
permissions, or contracts, because those areas tend to create cascading failures across the stack.

\subsection{Task Tracking and Communication}

Work was broken down by feature slices that mapped to ownership boundaries but were still traceable to
requirements. Typical slices included: API contract establishment, document persistence, realtime session
join, AI job orchestration, permission UI, and demo/deployment hardening.

Task tracking is best represented through GitHub issues or project-board items labeled by ownership area,
for example \texttt{backend}, \texttt{frontend}, \texttt{realtime}, \texttt{ai}, and \texttt{docs}. This makes
dependencies visible and helps the team see when one milestone is blocked by infrastructure, contracts,
or UI work owned elsewhere.

In practice, the team used GitHub for branch work, pull requests, issues, and durable code review
discussion, while Google Docs was used for collaborative drafting of the architecture and report content.
Lightweight team chat was used for short coordination loops, but technical decisions were expected to be
captured back into persistent artifacts such as ADRs, diagrams, PR discussions, README updates, and
contract notes instead of being lost in chat history.

\section{Development Methodology}

The most suitable methodology for this project is a lightweight Kanban/Scrum hybrid.

\begin{itemize}
  \item \textbf{Why not pure Scrum?} A student team with cross-cutting architecture tasks often has
  uneven availability and many infrastructure tasks that do not map cleanly to short user-story sprints.
  \item \textbf{Why not pure ad hoc Kanban?} The assignment has staged deliverables --- requirements,
  architecture, PoC, final MVP hardening --- that benefit from milestone-based planning.
\end{itemize}

In practice, the team can work in short milestone-oriented iterations while still prioritizing a Kanban
board by urgency and dependency. Backlog order should favor infrastructure and contract work that unblocks
multiple team members before polishing user-facing features.

Non-feature work such as writing tests, refining diagrams, splitting large files, or hardening local
startup scripts is treated as first-class work rather than invisible overhead. This is especially
important because the assignment explicitly evaluates repository organization, consistency, and the extent
to which the PoC matches the documented architecture.

\section{Risk Assessment}

\begin{longtable}{p{3.1cm}p{1.8cm}p{3.0cm}p{4.6cm}p{3.8cm}}
\caption{Technical risk assessment}\label{tab:risk-assessment}\\
\toprule
\textbf{Risk} & \textbf{Likelihood} & \textbf{Impact} & \textbf{Mitigation} & \textbf{Contingency} \\
\midrule
\endfirsthead
\toprule
\textbf{Risk} & \textbf{Likelihood} & \textbf{Impact} & \textbf{Mitigation} & \textbf{Contingency} \\
\midrule
\endhead
Realtime sync divergence or duplicate updates & Medium & Conflicting document views, lost trust in collaboration demo & Use Yjs instead of a custom algorithm, centralize realtime tests, preserve room state across reconnects, regression test save/reconnect paths & Fall back to manual save mode and present conflict-safe reload if realtime becomes unstable \\
AI provider latency, outage, or malformed responses & High & AI actions stall or fail during demo & Async job model, provider abstraction, deterministic stub fallback, timeout/error mapping tests & Switch demo provider to stub and keep editing workflow available \\
Backend and realtime configuration drift & Medium & Auth/session failures and invisible data mismatch across services & Shared root scripts, shared realtime secret, shared SQLite path, Docker Compose configuration & Reset with \texttt{npm run demo:reset} and verify health/docs endpoints before demo \\
Permission enforcement gaps between REST and realtime & Medium & Unauthorized edits or inconsistent collaboration behavior & Server-side RBAC on every REST route and websocket write path, viewer rejection tests, permission update events & Downgrade affected users to viewer mode and block live editing until policy is revalidated \\
Document-report divergence late in the semester & High & Lost marks despite working code & Keep architecture notes, OpenAPI, Mermaid sources, and final report in version control; explicitly document MVP deviations & Add a final alignment appendix and demo checklist so reviewers can map code to report claims quickly \\
\bottomrule
\end{longtable}


\section{Timeline and Milestones}

The project timeline must be described in terms of verifiable outcomes, not vague progress statements.
The milestones in Table~\ref{tab:milestones} reflect how the system could evolve through the semester
without losing alignment between requirements, architecture, and code.

\begin{table}[H]
\centering
\caption{Semester milestone plan}
\label{tab:milestones}
\begin{tabularx}{\textwidth}{p{2.1cm}p{4.2cm}X}
\toprule
\textbf{Milestone} & \textbf{Goal} & \textbf{Acceptance Criteria} \\
\midrule
M1 & Requirements and architecture baseline & Stakeholders, functional requirements, NFRs, user stories, C4 L1/L2/L3, data model, and ADR drafts reviewed by the team \\
M2 & Backend and API contract foundation & Auth, document create/load/save, permissions, sessions, AI job creation, and OpenAPI/contract notes verified by automated backend tests \\
M3 & Realtime collaboration slice & Two users can join the same document, see presence, and converge on edits through the realtime service with viewer write protection \\
M4 & AI and versioning integration & Async AI actions, AI policy controls, version history, revert, and export jobs demonstrated through frontend and backend integration tests \\
M5 & Submission hardening & Root scripts, Docker Compose, CI, final report, diagram sources, README, and video script completed and reviewed against the rubric \\
\bottomrule
\end{tabularx}
\end{table}

